\chapter{Research Gap and Problem Formulation}
\label{ch:gap}

\section{Identified Research Gap}
While significant progress has been made across temporal modeling, risk-sensitive rewards, and walk-forward validation protocols, existing studies have typically evaluated these mechanisms within isolated frameworks using divergent market universes, varied cost assumptions, and concurrent architectural changes. 

Based on the literature reviewed, there remains a clear need for controlled empirical comparisons that systematically isolate the contributions of temporal memory, risk-aware rewards, and turnover regularization under a common, realistic experimental setup. Without holding the asset universe, feature set, transaction-cost model, training budget, and chronological evaluation protocol constant, it is difficult to determine whether reported differences in out-of-sample performance stem from the sequence encoder, the reward function, turnover control, or differences in the experimental setup itself.

Table~\ref{tab:positioning} positions this thesis relative to the reviewed literature, highlighting the controlled single-factor ablation design adopted here.

\begin{table}[htbp]
\centering
\scriptsize
\setlength{\tabcolsep}{3.5pt}
\caption{Positioning of this work relative to the reviewed literature.}
\label{tab:positioning}
\begin{tabularx}{\textwidth}{>{\raggedright\arraybackslash}p{2.8cm}>{\centering\arraybackslash}X>{\centering\arraybackslash}p{3.1cm}>{\centering\arraybackslash}p{2.4cm}>{\centering\arraybackslash}p{2.5cm}>{\centering\arraybackslash}X}
\toprule
\textbf{Paper} & \textbf{Temporal Memory} & \textbf{Risk-Aware / Reward} & \textbf{Turnover Control} & \textbf{Continuous Allocation} & \textbf{Controlled Ablation} \\
\midrule
Millea (2021) \citep{millea2021deep} & & \cmark~\textit{(risk-aware rewards discussed)} & & & \\
Zou et al. (2024) \citep{zou2024novel} & \cmark & & & & \\
Huang et al. (2024) \citep{huang2024self} & & \cmark & & & \\
Liu et al. (2024) \citep{liu2024dynamic} & & & & \cmark~\textit{(portfolio weights)} & \\
Wang \& Liu (2025) \citep{wang2025risk} & & \cmark~\textit{(risk-sensitive reward)} & \cmark & \cmark & \\
Yang et al. (2022) \citep{yang2022selective} & & & & \cmark~\textit{(Dirichlet distribution)} & \\
\midrule
\textbf{This work} & \cmark & \cmark~\textit{(DSR)} & \cmark~\textit{(turnover penalty)} & \cmark & \cmark \\
\bottomrule
\end{tabularx}
\end{table}

\section{POMDP Formulation of Daily Portfolio Allocation}
\label{sec:pomdp_formulation}
We formulate daily portfolio allocation as a Partially Observable Markov Decision Process (POMDP), defined by the 7-tuple:
\begin{equation}
\mathcal{P} = \langle \mathcal{S}, \mathcal{A}, \mathcal{T}, \mathcal{R}, \Omega, \mathcal{O}, \gamma \rangle,
\end{equation}
where $\mathcal{S}$ is the unobservable state space, $\mathcal{A}$ is the action space, $\mathcal{T}$ represents transition dynamics, $\mathcal{R}$ is the reward function, $\Omega$ is the observation space, $\mathcal{O}$ is the observation probability distribution, and $\gamma \in [0, 1)$ is the discount factor.

\subsection{Why Daily Financial Trading is a POMDP}
In a fully observable Markov Decision Process (MDP), the Markov property requires that future state transitions depend solely on the current state and action: $\mathbb{P}(\mathbf{s}_{t+1} \mid \mathbf{s}_t, \mathbf{a}_t) = \mathbb{P}(\mathbf{s}_{t+1} \mid \mathbf{s}_t, \dots, \mathbf{s}_0)$. In financial trading, this assumption does not hold when applied directly to observable market indicators:
\begin{enumerate}
    \item \textbf{Unobservable True Market State $\mathcal{S}$:} The complete economic state of financial markets includes macroeconomic fundamentals, central bank expectations, institutional limit order book queues, private liquidity flows, and participant psychology. These variables are latent and cannot be directly observed by an agent trading from public daily price feeds.
    \item \textbf{Partial and Noisy Observations $\Omega$:} The agent receives only an observation vector $\mathbf{o}_t \in \Omega$ drawn from the observation model $\mathcal{O}(\mathbf{o}_t \mid \mathbf{s}_t)$, consisting of daily open, high, low, close prices, trading volumes, and derived technical indicators.
    \item \textbf{Non-Markovian Indicator Dynamics:} An identical single-day technical observation vector (e.g., RSI = 65, MACD crossover) can result in completely different future price movements depending on latent market conditions (such as the beginning of a sustained trend versus an institutional distribution phase).
\end{enumerate}
Because single-day observations $\mathbf{o}_t$ do not satisfy the Markov property, an agent benefits from maintaining an internal temporal representation or belief state $b(\mathbf{s}_t)$ over historical observation sequences $\mathbf{F}_t = [\mathbf{s}_{t-W+1}, \dots, \mathbf{s}_t]$. In this work, recurrent neural networks (LSTM) are used to construct this temporal summary $\mathbf{h}_t \approx b(\mathbf{s}_t)$ end-to-end, directly mitigating the effects of partial observability.

\subsection{Action Space $\mathcal{A}$ and Simplex Manifold}
The portfolio consists of $N = 10$ risky equities and $1$ cash account, yielding an action dimension of $M = N + 1 = 11$. The action vector $\mathbf{a}_t = [a_{t,1}, \dots, a_{t,N+1}]^\top$ represents target portfolio weights, strictly constrained to lie on the unit simplex:
\begin{equation}
\label{eq:simplex_constraint}
\mathcal{A} \equiv \Delta^N = \left\{ \mathbf{w}_t \in \mathbb{R}^{N+1} \;\middle|\; w_{i,t} \ge 0 \; \forall i, \quad \sum_{i=1}^{N} w_{i,t} + w_{\text{cash},t} = 1 \right\}.
\end{equation}
Following \citet{yang2022selective}, this constraint is enforced directly by parameterising the actor network over a multivariate Dirichlet distribution, eliminating the need for heuristic projection or unconstrained Gaussian clipping.
